\documentclass[10pt,a4paper]{amsart}
\usepackage[margin=19mm]{geometry}
\usepackage{amsmath,amssymb,mathtools}
\usepackage[scr=rsfs,cal=euler]{mathalfa}
\usepackage{xfrac}
\usepackage[unicode,hidelinks,bookmarks=false]{hyperref}
\newcommand{\ie}{i.e.\ }
\newcommand{\cf}{cf.\ }
\newcommand{\resp}{respectively\ }
\newcommand{\Subsection}{\subsection}
\providecommand{\nobreakdash}{\nobreak\hbox{-}}
\setlength{\emergencystretch}{2em}
\multlinegap0pt
\usepackage{amssymb}
\newtheorem{theorem}{Theorem}
\newcommand{\comment}[1]{\marginpar{\hfuzz=5ex
\scriptsize$\stackrel{?}{\longleftarrow}$\\\hspace*{-3ex}
{\fbox{\parbox{20ex}{\hfuzz=15ex \raggedright\sf#1 }}} }}
\title{Hodge decomposition for generalized Vekua spaces in higher dimensions}
\author{Briceyda B. Delgado}
\date{Source: arXiv:2305.15857v2 (CC BY 4.0)}
\begin{document}
\maketitle

\noindent\emph{Source and licence.} Hodge decomposition for generalized Vekua spaces in higher dimensions. Version arXiv:2305.15857v2 (CC BY 4.0), distributed by arXiv under the Creative Commons Attribution 4.0 International licence, http://creativecommons.org/licenses/by/4.0/, as recorded in the arXiv licence metadata for this exact version and shown on the abstract page https://arxiv.org/abs/2305.15857v2. Full licence notice: \texttt{licenses/arxiv-2305.15857-CC-BY-4.0.txt}.\par\smallskip

\noindent\textit{Editorial scope.} Counted unit: the article's theorem th:Hodge-decomposition (source0.tex lines 530--538). The article's complete original proof of it is delivered with it (source0.tex lines 539--566, one \texttt{proof} environment of 28 lines). No mathematics is written, repaired or re-derived: the statement and the proof are the article's own lines, in the article's own notation, and no retained line is substituted or re-flowed.  Retained uncounted as the source-local context the proof needs: lines 265--267; lines 451--453.  Nothing was omitted from any retained span, so the package carries no omission record.  No retained line contains a citation, so the wrapper carries no bibliography and no bibliography entry.  No retained line contains a figure environment, an image-inclusion command or drawing code, so no image is delivered and the package declares no asset dependency.  Editorial formatting only, in addition: an AMS \texttt{amsart} preamble that restates the article's own declarations and macros used by the retained lines, namely the environments \texttt{theorem} on separate counters, together with the standard packages those lines need.  The retained environments print in this document as Theorem 1 in order of appearance, whereas the article prints the same items with the numbering of its own full document.  The complete native source of the article as distributed in the arXiv e-print for this exact version is bundled unchanged as \texttt{sources/142016/source0.tex}.
\begin{align}\label{eq:Vekua_equation}
   D w=\alpha\, \overline{w}+\beta w, \quad \text{ in }\Omega.
\end{align}

\begin{equation}\label{eq:scalar-product}
   \langle u,v\rangle_{0,L^2}=\text{Sc }\int_{\Omega} \overline{u}v\, d\sigma, 
\end{equation}

\begin{theorem}\label{th:Hodge-decomposition}(\textbf{Hodge decomposition theorem})
%\comment{We need  the hypothesis of boundedness to have well defined the Neumann series}
Let $\Omega$ be a bounded Lipschitz domain, and let $\alpha, \beta\in L^{\infty}(\Omega,C\ell_{0,n})$.
%$\|\alpha\|_{\infty}+\|\beta\|_{\infty}\leq 1/\|T_{\Omega}\|$. 
Then the Hilbert space $L^2(\Omega,C\ell_{0,n})$ allows the orthogonal decomposition
\begin{align}\label{eq:Hodge-2}  L^2(\Omega,C\ell_{0,n})=A_{\alpha,\beta}^2(\Omega,C\ell_{0,n})\oplus (D-M^{\alpha}C-\overline{\beta}) W_0^{1,2}(\Omega,C\ell_{0,n}),
\end{align} 
under the scalar product \eqref{eq:scalar-product}.
\end{theorem}

\begin{proof}
It remains to prove that $(D-M^{\alpha}C-\overline{\beta}) W_0^{1,2}(\Omega,C\ell_{0,n})=A_{\alpha,\beta}^2(\Omega,C\ell_{0,n})^{\bot}$.
Let $h\in W_0^{1,2}(\Omega,C\ell_{0,n})$ and $w\in A_{\alpha,\beta}^2(\Omega,C\ell_{0,n})$. By the Gauss's theorem and from the fact that $h$ has vanishing boundary values, then
\begin{align}\label{eq:computations-orthogonal}
  \nonumber \langle w, (D-M^{\alpha}C-\overline{\beta})h \rangle_{0,L^2}&=\text{Sc }\int_{\Omega} \overline{w} \, (Dh-\overline{h}\alpha-\overline{\beta}h)\, d\sigma\\
  \nonumber &=\text{Sc }\int_{\Omega} -((\overline{w}D)h+\overline{w}\overline{h}\alpha+\overline{w}\overline{\beta}h)\, d\sigma+\text{Sc }\int_{\partial\Omega} \overline{w} \,\eta\,  h\, ds\\
    \nonumber &=\text{Sc }\int_{\Omega} -((\overline{w}D)h+w\overline{\alpha}h+\overline{w}\overline{\beta}h)\, d\sigma\\
    &=\text{Sc }\int_{\Omega} \overline{(D-\alpha C-\beta)w} \, h\, d\sigma,
\end{align}
where the third equality is justified by $\text{Sc }(\overline{w}\overline{h}\alpha)=\text{Sc }(w\overline{\alpha}h)$. Thus, 
\begin{equation*}%\label{eq:contention-1}
(D-M^{\alpha}C-\overline{\beta}) W_0^{1,2}(\Omega,C\ell_{0,n})\subset A_{\alpha,\beta}^2(\Omega,C\ell_{0,n})^{\bot}.
\end{equation*}
Now, we will prove that 
\begin{equation}\label{eq:contention-2}
(D-M^{\alpha}C-\overline{\beta})C_{c}^{\infty}(\Omega,C\ell_{0,n})^{\bot}\subset A_{\alpha,\beta}^2(\Omega,C\ell_{0,n}).
\end{equation}

Let $w \in (D-M^{\alpha}C-\overline{\beta})C_{c}^{\infty}(\Omega,C\ell_{0,n})^{\bot}$, then
\[
0=\nonumber \langle w, (D-M^{\alpha}C-\overline{\beta})g \rangle_{0,L^2},
\]
for all $g\in C_{c}^{\infty}(\Omega,C\ell_{0,n})$. Using again \eqref{eq:computations-orthogonal}, we obtain
\begin{align*}%\label{eq:converse-1}
0=\text{Sc }\int_{\Omega} \overline{(D-\alpha C-\beta)w} \, g\, d\sigma,
\end{align*}
for all $g\in C_{c}^{\infty}(\Omega,C\ell_{0,n})$. That is, $w$ is a weak solution of the Vekua equation \eqref{eq:Vekua_equation}, which implies \eqref{eq:contention-2}. Applying the fact that $\overline{C_{c}^{\infty}(\Omega,C\ell_{0,n})}=W^{1,2}(\Omega,C\ell_{0,n})$ in \eqref{eq:contention-2} we readily obtain that $A_{\alpha,\beta}^2(\Omega,C\ell_{0,n})^{\bot}\subset (D-M^{\alpha}C-\overline{\beta}) W_0^{1,2}(\Omega,C\ell_{0,n})$, which completes the proof.
\end{proof}

\end{document}
